\documentclass[10pt,a4paper]{amsart}
\usepackage[margin=19mm]{geometry}
\usepackage{amsmath,amssymb,mathtools}
\usepackage[scr=rsfs,cal=euler]{mathalfa}
\usepackage{xfrac}
\usepackage[unicode,hidelinks,bookmarks=false]{hyperref}
\newcommand{\ie}{i.e.\ }
\newcommand{\cf}{cf.\ }
\newcommand{\resp}{respectively\ }
\newcommand{\Subsection}{\subsection}
\providecommand{\nobreakdash}{\nobreak\hbox{-}}
\setlength{\emergencystretch}{2em}
\multlinegap0pt
\usepackage{amssymb}
\usepackage{mathtools}
\usepackage{tikz}
\newtheorem{lemma}{Lemma}
\newtheorem{proposition}{Proposition}
\newtheorem{theorem}{Theorem}
\newcommand{\Ee}{\mathbb E}
\newcommand{\Pp}{\mathbb P}
\newcommand{\R}{\mathbb R}
\newcommand{\TV}{\mathrm{TV}}
\newcommand{\Z}{\mathbb Z}
\newcommand{\abs}[1]{\left|#1\right|}
\newcommand{\cD}{\mathcal D}
\newcommand{\cF}{\mathcal F}
\newcommand{\cH}{\mathcal H}
\newcommand{\ceil}[1]{\left\lceil #1\right\rceil}
\newcommand{\floor}[1]{\left\lfloor #1\right\rfloor}
\newcommand{\id}{\mathrm{id}}
\newcommand{\modn}[1]{\overline{#1}}
\newcommand{\norm}[1]{\left\lVert #1\right\rVert}
\newcommand{\one}{\mathbf 1}
\newcommand{\pospart}[1]{\left(#1\right)_{+}}
\title{Periodic ASEP and random walks on affine Hecke Algebra}
\author{Alexey Bufetov, Nimisha Pahuja}
\date{Source: arXiv:2609.10358v1 (CC BY 4.0)}
\begin{document}
\maketitle

\noindent\emph{Source and licence.} Periodic ASEP and random walks on affine Hecke Algebra. Version arXiv:2609.10358v1 (CC BY 4.0), distributed by arXiv under the Creative Commons Attribution 4.0 International licence, http://creativecommons.org/licenses/by/4.0/, as recorded in the arXiv licence metadata for this exact version and shown on the abstract page https://arxiv.org/abs/2609.10358v1. Full licence notice: \texttt{licenses/arxiv-2609.10358-CC-BY-4.0.txt}.\par\smallskip

\noindent\textit{Editorial scope.} Counted unit: the article's theorem thm:limit-shape-step (source0.tex lines 855--928). The article's complete original proof of it is delivered with it (source0.tex lines 930--1107, one \texttt{proof} environment of 178 lines). No mathematics is written, repaired or re-derived: the statement and the proof are the article's own lines, in the article's own notation, and no retained line is substituted or re-flowed.  Retained uncounted as the source-local context the proof needs: lines 477--484; lines 505--519; lines 598--603; lines 606--619; lines 797--804.  Nothing was omitted from any retained span, so the package carries no omission record.  No retained line contains a citation, so the wrapper carries no bibliography and no bibliography entry.  No retained line contains a figure environment, an image-inclusion command or drawing code, so no image is delivered and the package declares no asset dependency.  Editorial formatting only, in addition: an AMS \texttt{amsart} preamble that restates the article's own declarations and macros used by the retained lines, namely the environments \texttt{lemma}, \texttt{proposition}, \texttt{theorem} on separate counters, together with the standard packages those lines need.  The retained environments print in this document as Lemma 1, Proposition 1, Theorem 1 in order of appearance, whereas the article prints the same items with the numbering of its own full document.  The complete native source of the article as distributed in the arXiv e-print for this exact version is bundled unchanged as \texttt{sources/209902/source0.tex}.
	\begin{lemma}\label{lem:ring-mixing}
		For every $\varepsilon>0$, there exists $T_\varepsilon<\infty$ such that
		\begin{equation}\label{eq:ring-mixing}
			\sup_{\omega\in S_n}
			\norm{P_t(\omega,\cdot)-\pi_n^q}_{\TV}<\varepsilon,
			\qquad t\ge T_\varepsilon.
		\end{equation}
	\end{lemma}

	\begin{lemma}\label{lem:ring-velocity}
		For every initial ring state and every $i\in[n]$,
		\begin{equation}\label{eq:ring-velocity-lln}
			\frac{X_i(t)}{t}\longrightarrow v_i
			\qquad\text{almost surely and in }L^1,
		\end{equation}
		where
		\begin{equation}\label{eq:ring-velocities}
			v_i=\frac{(1-q)(n+1-2i)}{n-1}.
		\end{equation}
		In particular,
		\begin{equation}\label{eq:strict-velocity-order}
			v_1>v_2>\cdots>v_n.
		\end{equation}
	\end{lemma}

	\begin{equation}\label{eq:safe-ring-word}
		\Omega(w)(r)=k
		\quad\Longleftrightarrow\quad
		Y_{a_k}\equiv r\pmod n,
		\qquad r,k\in[n].
	\end{equation}

	\begin{lemma}[Safe-region coupling]\label{lem:safe-region-coupling}
		Start the affine walk from $w\in\cD$ and a ring ASEP from $\Omega(w)$, using
		the same marked clocks.  Let $X_k(t)$ be the lifted displacement of ring
		particle $k$, with $X_k(0)=0$.  Up to the first time at which the affine walk
		leaves $\cD$,
		\begin{equation}\label{eq:safe-coupling}
			Y_{a_k}(t)=Y_{a_k}(0)+X_k(t),
			\qquad k\in[n].
		\end{equation}
		Moreover, if the $n$ paths on the right-hand side of
		\eqref{eq:safe-coupling} remain pairwise at distance greater than $n$ for all
		$t\ge0$, then the affine walk never leaves $\cD$ and
		\eqref{eq:safe-coupling} holds for all $t\ge0$.
	\end{lemma}

	\begin{proposition}\label{prop:eventual-separation}
		From every initial affine permutation, and in particular from the identity,
		there exists almost surely a finite random time $T_\mathrm{sep}$ such that
		\begin{equation}\label{eq:eventual-safe-region}
			w_t\in\cD,
			\qquad \mbox{for all $\ t \ge T_\mathrm{sep}$ }.
		\end{equation}
	\end{proposition}

	\begin{theorem}[Limit shape and local equilibrium for step initial data]
		\label{thm:limit-shape-step}
		Let the affine Hecke walk start from $w_0=\id$, and let
		\begin{equation}\label{eq:chi-from-wt}
			\chi_t(x)=\one_{\{w_t(x)\le0\}}
		\end{equation}
		be its one-species projection.  Recall the leader velocities
		\begin{equation}\label{eq:theorem-velocities}
			v_k=\frac{(1-q)(n+1-2k)}{n-1},
			\qquad k\in[n].
		\end{equation}
		Then the following assertions hold.
		
		\smallskip
		\noindent\emph{(i) Leader law of large numbers.}
		Almost surely, there is a random permutation $(a_1,\ldots,a_n)$ such that,
		for all sufficiently large $t$,
		\begin{equation}\label{eq:eventual-leader-order}
			Y_{a_1}(t)>Y_{a_2}(t)>\cdots>Y_{a_n}(t),
		\end{equation}
		and
		\begin{equation}\label{eq:leader-velocity-limit}
			\frac{Y_{a_k}(t)}{t}\longrightarrow v_k,
			\qquad k\in[n].
		\end{equation}
		
		\smallskip
		\noindent\emph{(ii) Integrated limit shape.}
		Let
		\begin{equation}\label{eq:particle-count}
			N_t(x)=\sum_{y>x}\chi_t(y),
			\qquad x\in\Z.
		\end{equation}
		Then, almost surely,
		\begin{equation}\label{eq:uniform-height-limit}
			\sup_{u\in\R}
			\abs{\frac{N_t(\floor{ut})}{t}-\cH(u)}
			\longrightarrow0,
		\end{equation}
		where
		\begin{equation}\label{eq:limit-height-function}
			\cH(u)=\frac1n\sum_{k=1}^n\pospart{v_k-u}.
		\end{equation}
		In particular, away from the break points $v_1,\ldots,v_n$, the macroscopic
		density $\rho(u)=-\cH'(u)$ is
		\begin{equation}\label{eq:density-profile}
			\rho(u)=
			\begin{cases}
				1,&u<v_n,\\[1mm]
				\dfrac{K}{n},&v_{K+1}<u<v_K,
				\quad 1\le K\le n-1,\\[2mm]
				0,&u>v_1.
			\end{cases}
		\end{equation}
		
		\smallskip
		\noindent\emph{(iii) Local equilibrium.}
		Let $x_t\in\Z$ be deterministic and satisfy $x_t/t\to u$, where
		$u\notin\{v_1,\ldots,v_n\}$.  Let $K=K(u)$ be the unique integer in
		$\{0,\ldots,n\}$ such that $\rho(u)=K/n$.  Then, for every fixed
		$N\ge1$,
		\begin{equation}\label{eq:local-equilibrium}
			\bigl(\chi_t(x_t+j)\bigr)_{j=0}^{N-1}
			\ \Longrightarrow\ 
			\bigl(\eta(j)\bigr)_{j=0}^{N-1},
			\qquad \eta\sim\nu_{n,K}.
		\end{equation}
		Equivalently, for every bounded function $f:\{0,1\}^N\to\R$,
		\begin{equation}\label{eq:local-equilibrium-test-functions}
			\lim_{t\to\infty}
			\Ee\left[f\bigl((\chi_t(x_t+j))_{j=0}^{N-1}\bigr)\right]
			=\int f\bigl((\eta(j))_{j=0}^{N-1}\bigr)\,\nu_{n,K}(d\eta).
		\end{equation}
	\end{theorem}

	\begin{proof}
		We prove the three parts in turn.
		
		\medskip
		\noindent\emph{Step 1: asymptotic velocities of the leaders.}
		For each deterministic integer $T\ge0$, construct an auxiliary ring ASEP
		from the graphical marks after time $T$.  If $w_T\in\cD$, order the leaders
		from right to left, form the ring word $\Omega(w_T)$ as in
		\eqref{eq:safe-ring-word}, and start the ring process from this state.  If
		$w_T\notin\cD$, start it from an arbitrary fixed state.  Denote the lifted
		displacement of ring particle $k$ in this auxiliary process by
		$X_k^{(T)}(s)$.
		
		Conditionally on $\cF_T$, the initial ring state is fixed and the future
		marks are independent of $\cF_T$.  Lemma~\ref{lem:ring-velocity} therefore
		implies, with conditional probability one,
		\[
		\frac{X_k^{(T)}(s)}s\longrightarrow v_k,
		\qquad k\in[n].
		\]
		For each fixed $T$ this holds on an event of probability one; intersecting
		these events over $T\in\Z_{\ge0}$ preserves probability one.
		
		By Proposition~\ref{prop:eventual-separation}, almost surely there is an
		integer $T$ such that $w_t\in\cD$ for every $t\ge T$.  Fix such a $T$ on the
		full-probability event just constructed, and write $a_1,\ldots,a_n$ for the
		right-to-left order of the leaders at time $T$.  This order cannot change
		while the walk remains in $\cD$.  Lemma~\ref{lem:safe-region-coupling}
		gives, for every $s\ge0$,
		\begin{equation}\label{eq:theorem-safe-coupling}
			Y_{a_k}(T+s)-Y_{a_k}(T)=X_k^{(T)}(s),
			\qquad k\in[n].
		\end{equation}
		Dividing by $T+s$ and letting $s\to\infty$ proves
		\eqref{eq:leader-velocity-limit}, and hence part~(i).
		
		\medskip
		\noindent\emph{Step 2: the integrated limit shape.}
		For each $i\in[n]$, the particles whose labels are congruent to $i$ modulo
		$n$ have labels and positions
		\begin{equation}\label{eq:class-particles}
			i+nk
			\quad\text{and}\quad
			Y_i(t)+nk,
			\qquad k\in\Z,
		\end{equation}
		respectively.  Since $i+nk\le0$ if and only if $k\le-1$, the contribution of
		this residue class to $N_t(x)$ is
		\begin{equation}\label{eq:class-count-exact}
			\#\{k\le-1:Y_i(t)+nk>x\}
			=\pospart{\ceil{\frac{Y_i(t)-x}{n}}-1}.
		\end{equation}
		For every $z\in\R$,
		\[
		\abs{\pospart{\ceil{z/n}-1}-\pospart{z/n}}\le1.
		\]
		Summing over the $n$ residue classes therefore gives
		\begin{equation}\label{eq:height-approx-leaders}
			\abs{N_t(x)-\frac1n\sum_{i=1}^n\pospart{Y_i(t)-x}}
			\le n.
		\end{equation}
		Since $(a-u)_+$ is $1$-Lipschitz in each variable,
		\eqref{eq:height-approx-leaders}, part~(i), and
		$\abs{\floor{ut}/t-u}\le t^{-1}$ yield
		\begin{align*}
			&\sup_{u\in\R}
			\abs{\frac{N_t(\floor{ut})}{t}
				-\frac1n\sum_{k=1}^n\pospart{v_k-u}}\\
			&\qquad\le \frac nt+\frac1t
			+\max_{1\le k\le n}
			\abs{\frac{Y_{a_k}(t)}t-v_k}
			\longrightarrow0.
		\end{align*}
		This proves \eqref{eq:uniform-height-limit}.  Formula
		\eqref{eq:density-profile} follows by differentiating
		\eqref{eq:limit-height-function} away from its break points.
		
		\medskip
		\noindent\emph{Step 3: local equilibrium.}
		For a deterministic integer $T\ge0$, set
		\begin{equation}\label{eq:AT-definition}
			A_T=\{w_s\in\cD\text{ for every }s\ge T\}.
		\end{equation}
		The events $A_T$ increase with $T$, and Proposition
		\ref{prop:eventual-separation} implies
		\begin{equation}\label{eq:AT-probability}
			\Pp(A_T)\longrightarrow1
			\qquad(T\to\infty).
		\end{equation}
		Let $\omega^{(T)}$ be the auxiliary ring process from Step~1 and define its
		$K$-particle projection by
		\begin{equation}\label{eq:ring-K-projection}
			\zeta_s^{(T,K)}(r)
			=\one_{\{\omega_s^{(T)}(r)\le K\}},
			\qquad r\in\Z/n\Z.
		\end{equation}
		
		On $A_T$, let $a_1,\ldots,a_n$ be the right-to-left order of the leaders at
		time $T$.  The safe-region coupling implies
		\begin{equation}\label{eq:ring-label-at-leader-residue}
			\omega_{t-T}^{(T)}\bigl(\modn{Y_{a_k}(t)}\bigr)=k,
			\qquad k\in[n],\quad t\ge T.
		\end{equation}
		Indeed, ring particle $k$ starts at residue
		$\modn{Y_{a_k}(T)}$ and its lifted displacement is
		$Y_{a_k}(t)-Y_{a_k}(T)$.
		
		Fix $y\in\Z$, and let $k$ be the unique index for which
		$y\equiv Y_{a_k}(t)\pmod n$.  For a unique $m\in\Z$,
		\[
		y=Y_{a_k}(t)+nm,
		\qquad
		w_t(y)=a_k+nm.
		\]
		As $a_k\in[n]$, it follows that
		\begin{equation}\label{eq:occupation-via-leader}
			\chi_t(y)=1
			\quad\Longleftrightarrow\quad
			m\le-1
			\quad\Longleftrightarrow\quad
			Y_{a_k}(t)\ge y+n.
		\end{equation}
		Use the conventions $v_0=+\infty$ and $v_{n+1}=-\infty$.  By the definition
		of $K$,
		\[
		v_K>u>v_{K+1}.
		\]
		Part~(i), $x_t/t\to u$, and the fact that $N$ is fixed imply that, almost
		surely on $A_T$, for all sufficiently large $t$, simultaneously for
		$0\le j<N$,
		\begin{equation}\label{eq:block-coupling}
			\chi_t(x_t+j)
			=\zeta_{t-T}^{(T,K)}\bigl(\modn{x_t+j}\bigr).
		\end{equation}
		
		It remains to pass from this pathwise comparison to the unconditional
		limiting law. Set
		\[
		d(s)=\sup_{\omega\in S_n}
		\norm{P_s(\omega,\cdot)-\pi_n^q}_{\TV}.
		\]
		By Lemma~\ref{lem:ring-mixing}, $d(s)\to0$ as $s\to\infty$.
		The projection $\omega\mapsto(\one_{\{\omega(r)\le K\}})_{r\in\Z/n\Z}$
		sends $\pi_n^q$ to the uniform law on $K$-particle ring configurations.
		
		For fixed $T$, define the two blocks
		\[
		B_t=\bigl(\chi_t(x_t+j)\bigr)_{j=0}^{N-1},\qquad
		B_t^{(T)}=
		\bigl(\zeta_{t-T}^{(T,K)}(\modn{x_t+j})\bigr)_{j=0}^{N-1},
		\qquad t\ge T.
		\]
		The initial state of $\omega^{(T)}$ is $\cF_T$-measurable, and its driving
		marks after time $T$ are independent of $\cF_T$. Thus its unconditional
		law at time $t-T$ is within $d(t-T)$ of $\pi_n^q$ in total variation.
		For a bounded function $f:\{0,1\}^N\to\R$, write
		\[
		c_f=\int f\bigl((\eta(j))_{j=0}^{N-1}\bigr)\,\nu_{n,K}(d\eta).
		\]
		Rotation invariance of the uniform ring law handles the possibly varying
		residue of $x_t$, and gives
		\[
		\left|\Ee[f(B_t^{(T)})]-c_f\right|
		\le 2\norm{f}_\infty d(t-T).
		\]
		On $A_T$, equation~\eqref{eq:block-coupling} holds for all sufficiently
		large $t$, almost surely. Dominated convergence therefore implies
		\[
		\limsup_{t\to\infty}\Pp(B_t\ne B_t^{(T)})\le\Pp(A_T^c).
		\]
		Combining these estimates, for every fixed deterministic integer $T$,
		\[
		\limsup_{t\to\infty}\left|\Ee[f(B_t)]-c_f\right|
		\le2\norm{f}_\infty\Pp(A_T^c).
		\]
		Finally let $T\to\infty$ and use \eqref{eq:AT-probability}. This proves
		\eqref{eq:local-equilibrium-test-functions}, and hence part~(iii).
	\end{proof}

\end{document}
